\section{功能性需求}

\subsection{功能模块总览}

\srsfigure{figures/srs/09-feature-map.pdf}{系统功能模块图}{fig:feature-map}

系统功能分为公共账号与通知、顾客端、商家端、骑手端、管理端、智能服务、营销资产与个性化七个模块。每项需求使用稳定编号，后续接口、数据库、测试和验收材料必须引用同一编号。

\subsection{系统总用例}

\srsfigure{figures/srs/10-system-use-case.pdf}{系统总用例图}{fig:system-use-case}

总用例图只展示四类业务角色、游客和外部服务与一级用例的关系，各角色内部用例见后续分图和用例表。

\subsection{用例编写规则}

每个用例包含执行者、触发条件、前置条件、基本流程、异常与边界、后置条件、业务规则和可直接执行的验收标准。需求中的“顺序”、“仅本人”、“不重复”等限制均应由后端再次校验，不得只依赖前端隐藏按钮。
